\documentclass[aspectratio=169, 10pt]{beamer}

% --- Theme Setup ---
\usetheme{Madrid}
\usecolortheme{beaver} % A red theme (often distinctive for Econ)
\setbeamertemplate{navigation symbols}{} 

% --- Packages ---
\usepackage{booktabs} 
\usepackage{graphicx}
\usepackage{amsmath}
\usepackage{xcolor}
\usepackage{tcolorbox}

% --- Metadata ---
\title{Unit 2: Economic Indicators \& The Business Cycle}
\subtitle{Measurement of Economic Performance}
\author{Doğukan Günkördü}
\date{\today}

\begin{document}

% --- Title ---
\begin{frame}
    \titlepage
\end{frame}

\begin{frame}{Unit Roadmap}
    \tableofcontents
\end{frame}

% =========================================================
% SECTION 1: GDP
% =========================================================
\section{Gross Domestic Product (GDP)}

\begin{frame}{1.1 The Circular Flow Matrix}
    The economy consists of two key actors and two key markets:
    \begin{columns}
        \column{0.5\textwidth}
        \begin{itemize}
            \item \textbf{Households:} Supply resources (labor) and demand goods.
            \item \textbf{Firms:} Demand resources and supply goods.
            \item \textbf{Product Market:} Where goods/services are sold.
            \item \textbf{Resource (Factor) Market:} Where land, labor, and capital are sold.
        \end{itemize}
        
        \column{0.5\textwidth}
        \begin{alertblock}{Key Insight}
        One person's spending is another person's income.
        \end{alertblock}
    \end{columns}
\end{frame}

\begin{frame}{1.2 Gross Domestic Product (GDP)}
    \textbf{Definition:} The dollar value of all \underline{final} goods and services produced within a country's borders in one year.
    
    \vspace{0.5cm}
    \textbf{The Three Rules of GDP:}
    \begin{enumerate}
        \item \textbf{Excludes Financial Transactions:} Stocks, bonds, and Real Estate transfers are NOT production.
        \item \textbf{Excludes Transfer Payments:} Welfare, Social Security, Stimulus checks (no production occurred).
        \item \textbf{Excludes Second-Hand Sales:} Used cars, thrift store clothes (counted when they were new).
    \end{enumerate}
\end{frame}

\begin{frame}{1.3 Components of GDP (Expenditures Approach)}
    $$ \text{GDP} = C + I + G + X_n $$
    
    \begin{description}
        \item[Consumption (C):] Spending by households. (Durable vs Non-durable goods). \textit{~70\% of US GDP.}
        \item[Investment (I):] Spending by businesses on capital (machines), new construction (houses), and \textbf{changes in inventory}.
        \item[Government (G):] Goods and services (schools, tanks). \textbf{NOT} transfer payments.
        \item[Net Exports ($X_n$):] Exports ($X$) - Imports ($M$).
    \end{description}
\end{frame}

\begin{frame}{Check for Understanding: GDP}
    \begin{exampleblock}{Question}
    Which of the following is included in US GDP?
    \end{exampleblock}
    
    \begin{enumerate}[A]
        \item A monthly Social Security check received by a retiree.
        \item The purchase of a used Toyota on Craigslist.
        \item The purchase of 100 shares of Google stock.
        \item A haircut bought at a local salon.
        \item Tires bought by Ford to put on a new truck.
    \end{enumerate}
    \pause
    \textbf{Answer:} D (Service). \textit{A, C are transfers/financial. B is used. E is an intermediate good.}
\end{frame}

% =========================================================
% SECTION 2: INFLATION
% =========================================================
\section{Inflation and Price Indices}

\begin{frame}{2.1 Nominal vs. Real GDP}
    \begin{itemize}
        \item \textbf{Nominal GDP:} GDP measured in current prices. It does not account for inflation.
        \item \textbf{Real GDP:} GDP expressed in constant, or unchanging, dollars. Adjusted for inflation.
    \end{itemize}
    
    \begin{block}{The GDP Deflator}
    A measure of the price level calculated as the ratio of nominal GDP to real GDP.
    $$ \text{GDP Deflator} = \frac{\text{Nominal GDP}}{\text{Real GDP}} \times 100 $$
    \end{block}
\end{frame}

\begin{frame}{2.2 Measuring Inflation: The CPI}
    \textbf{Consumer Price Index (CPI):} Measures the price of a standard "market basket" of goods purchased by a typical urban consumer.
    
    $$ \text{CPI} = \frac{\text{Price of Market Basket in Current Year}}{\text{Price of Market Basket in Base Year}} \times 100 $$
    
    \vspace{0.3cm}
    \textbf{The Inflation Rate Formula:}
    $$ \text{\% Change} = \frac{\text{Year}_2 - \text{Year}_1}{\text{Year}_1} \times 100 $$
\end{frame}

\begin{frame}{2.3 Winners and Losers of Inflation}
    \textbf{Unexpected} inflation redistributes wealth.
    
    \begin{columns}
        \column{0.5\textwidth}
        \begin{tcolorbox}[colback=green!10, title=Who is HELPED?]
        \begin{itemize}
            \item \textbf{Borrowers:} They repay loans with dollars that have lower purchasing power.
            \item A business where the price of the product increases faster than the price of resources.
        \end{itemize}
        \end{tcolorbox}
        
        \column{0.5\textwidth}
        \begin{tcolorbox}[colback=red!10, title=Who is HURT?]
        \begin{itemize}
            \item \textbf{Lenders:} They are paid back with less valuable dollars.
            \item \textbf{Savers:} Money in the bank loses real value.
            \item People on \textbf{Fixed Incomes}.
        \end{itemize}
        \end{tcolorbox}
    \end{columns}
\end{frame}

% =========================================================
% SECTION 3: UNEMPLOYMENT
% =========================================================
\section{Unemployment}

\begin{frame}{3.1 Defining the Labor Force}
    \textbf{Population Division:}
    \begin{enumerate}
        \item Under 16 or Institutionalized (Prison/Military).
        \item \textbf{Not in Labor Force:} Retirees, full-time students, stay-at-home parents, \textbf{discouraged workers}.
        \item \textbf{Labor Force:} 16+, willing/able to work.
    \end{enumerate}
    
    \vspace{0.5cm}
    $$ \text{Unemployment Rate} = \frac{\text{\# Unemployed}}{\text{\# in Labor Force}} \times 100 $$
\end{frame}

\begin{frame}{3.2 Three Types of Unemployment}
    \begin{block}{1. Frictional (Good/Normal)}
    "Between jobs", voluntary, qualified workers with transferable skills. E.g., College grad looking for first job.
    \end{block}
    
    \begin{block}{2. Structural (Creative Destruction)}
    Changes in the labor force make some skills obsolete. The jobs will not come back. E.g., VCR repairman replaced by streaming.
    \end{block}
    
    \begin{block}{3. Cyclical (Bad)}
    Caused by a recession. Demand for goods falls $\rightarrow$ demand for labor falls.
    \end{block}
\end{frame}

\begin{frame}{3.3 The Natural Rate of Unemployment (NRU)}
    \begin{itemize}
        \item The economy is at \textbf{Full Employment} when there is \textbf{no Cyclical Unemployment}.
        \item Full employment does NOT mean 0\% unemployment.
    \end{itemize}
    
    $$ \text{NRU} = \text{Frictional} + \text{Structural} $$
    
    \vspace{0.5cm}
    \textit{In the US, the NRU is typically around 4-6\%.}
\end{frame}

\begin{frame}{3.4 Criticisms of the Unemployment Rate}
    The official rate (U3) can be misleading because of:
    
    \begin{itemize}
        \item \textbf{Discouraged Workers:} People who have given up looking for a job. They are no longer in the labor force, so the unemployment rate \textit{drops} even though the situation is worse.
        \item \textbf{Underemployed Workers:} Part-time workers who want full-time work are counted as "employed."
    \end{itemize}
\end{frame}

% =========================================================
% SECTION 4: THE BUSINESS CYCLE
% =========================================================
\section{The Business Cycle}

\begin{frame}{4.1 Phases of the Business Cycle}
    The US economy grows over time, but it is not linear.
    
    \begin{itemize}
        \item \textbf{Peak:} Temporary maximum. Low unemployment, high inflation.
        \item \textbf{Recession (Contraction):} Decline in real GDP for 2 consecutive quarters (6 months). Unemployment rises.
        \item \textbf{Trough:} The bottom. High unemployment.
        \item \textbf{Expansion (Recovery):} Real GDP rises. Inflation rises.
    \end{itemize}
\end{frame}

\begin{frame}{4.2 Output Gaps}
    \begin{itemize}
        \item \textbf{Potential Output ($Y_P$):} The level of Real GDP produced when the economy is at the Natural Rate of Unemployment (Full Employment).
        \item \textbf{Actual Output ($Y$):} What we are actually producing.
    \end{itemize}
    
    \vspace{0.5cm}
    \begin{alertblock}{The Gaps}
    If $Y < Y_P$: \textbf{Recessionary Gap} (High Unemployment).\\
    If $Y > Y_P$: \textbf{Inflationary Gap} (High Inflation, Overheating).
    \end{alertblock}
\end{frame}

% =========================================================
% PRACTICE SECTION
% =========================================================
\section{Practice \& Application}

\begin{frame}{Practice Calculation: CPI}
    \textbf{Basket:} 10 Gallons of Milk + 5 Loaves of Bread.
    \begin{itemize}
        \item \textbf{Year 1 (Base):} Milk \$2, Bread \$1.
        \item \textbf{Year 2:} Milk \$2.50, Bread \$2.
    \end{itemize}
    
    \vspace{0.2cm}
    \textbf{Calculate Cost of Basket:}
    \begin{itemize}
        \item Year 1: $(10 \times 2) + (5 \times 1) = \$25$
        \item Year 2: $(10 \times 2.5) + (5 \times 2) = \$35$
    \end{itemize}
    
    \vspace{0.2cm}
    \textbf{Calculate CPI for Year 2:}
    $$ \frac{35}{25} \times 100 = 140 $$
    \textit{Inflation is 40\% from the base year.}
\end{frame}

\begin{frame}{Summary: The Big 3 Indicators}
    \begin{table}[]
    \centering
    \begin{tabular}{|l|l|l|}
    \hline
    \textbf{Indicator} & \textbf{What it measures} & \textbf{Goal} \\ \hline
    \textbf{Real GDP} & Economic Growth & $\approx 2-3\%$ growth \\ \hline
    \textbf{Unemployment} & Labor utilization & $\approx 4-6\%$ (NRU) \\ \hline
    \textbf{Inflation} & Change in price level & $\approx 2\%$ \\ \hline
    \end{tabular}
    \end{table}
\end{frame}

\begin{frame}
    \centering
    \Huge \textbf{Questions?}
\end{frame}

\end{document}
